%% =============================================================
%%  Otimização Estocástica Bancária do NIM para Guidance
%%  ALM-ALCO sob Contornos Regulatórios em Mercados Completos
%%  Relatório Técnico --- Doutorado PPGEE Poli-USP
%%
%%  Compilar com:
%%    pdflatex alm_bank_relatorio.tex
%%    pdflatex alm_bank_relatorio.tex   (2a passagem: TOC, refs)
%%
%%  Dependências (TeX Live / MiKTeX --- download automático):
%%    lmodern, microtype, amsmath, amssymb, bm,
%%    geometry, setspace, graphicx, float, booktabs,
%%    array, multirow, tabularx, xcolor, colortbl,
%%    hyperref, natbib, fancyhdr, listings, tcolorbox,
%%    babel (brazil), inputenc, fontenc, caption
%% =============================================================

\documentclass[12pt, a4paper, oneside]{article}

%% --- Encoding e idioma --------------------------------------
\usepackage[utf8]{inputenc}
\usepackage[T1]{fontenc}
\usepackage[brazil]{babel}

%% --- Fonte --------------------------------------------------
\usepackage{lmodern}
\usepackage{microtype}

%% --- Matematica ---------------------------------------------
\usepackage{amsmath, amssymb, amsthm}
\usepackage{bm}

%% --- Layout -------------------------------------------------
\usepackage[
  top=3.0cm, bottom=2.0cm,
  left=3.0cm, right=2.0cm
]{geometry}
\usepackage{setspace}
\onehalfspacing

%% --- Figuras e tabelas --------------------------------------
\usepackage{graphicx}
\usepackage{float}
\usepackage{booktabs}
\usepackage{array}
\usepackage{multirow}
\usepackage{tabularx}
\usepackage{xcolor}
\usepackage{colortbl}
\usepackage{caption}   % permite \captionsetup e legenda customizada

%% --- Cores --------------------------------------------------
\definecolor{almblue}{RGB}{21, 101, 192}
\definecolor{almgreen}{RGB}{0, 140, 90}
\definecolor{almred}{RGB}{183, 28, 28}
\definecolor{almgray}{RGB}{240, 243, 247}
\definecolor{almgraytext}{RGB}{80, 100, 120}

%% --- Hyperlinks (deve vir ANTES de natbib) ------------------
\usepackage[
  colorlinks  = true,
  linkcolor   = almblue,
  citecolor   = almblue,
  urlcolor    = almblue,
  anchorcolor = almblue,
  pdftitle    = {Otimizacao Estocástica Bancária do NIM para Guidance
                 ALM-ALCO sob Contornos Regulatórios em Mercados Completos},
  pdfauthor   = {Walter Correa Neto},
  pdfsubject  = {Relatório Técnico de Doutorado PPGEE Poli-USP},
  pdfkeywords = {ALM; ALCO; NIM; otimizacao estocástica; IRRBB; LCR; NSFR; RAROC},
]{hyperref}

%% --- natbib: citações autor-data, links via hyperref --------
%%     Funciona diretamente com \begin{thebibliography} manual.
%%     \cite{key}      → (AUTOR, ANO)   clicável
%%     \citet{key}     → AUTOR (ANO)    clicável
%%     Cada link leva à entrada correspondente nas referências.
\usepackage[round, colon, authoryear, sort]{natbib}
\setcitestyle{aysep={,}}   % separador autor-ano: vírgula

%% --- Fonte e nota ABNT para tabelas e figuras ---------------
%%     \fonte{texto}  → "Fonte: texto"  abaixo do float
%%     \nota{texto}   → "Nota: texto"   abaixo do float
%%     Posicionados como parágrafo dentro do ambiente figure/table,
%%     APÓS o \caption, usando \par\smallskip com fonte menor.
\newcommand{\fonte}[1]{%
  \par\smallskip
  \noindent{\footnotesize\textbf{Fonte:} #1}%
}
\newcommand{\nota}[1]{%
  \par\smallskip
  \noindent{\footnotesize\textbf{Nota:} #1}%
}

%% --- Cabeçalho e rodapé -------------------------------------
\usepackage{fancyhdr}
\pagestyle{fancy}
\fancyhf{}
\fancyhead[L]{\small\textcolor{almgraytext}{%
  Otimização Estocástica Bancária do NIM --- Doutorado PPGEE Poli-USP}}
\fancyhead[R]{\small\textcolor{almgraytext}{W.\ Correa Neto}}
\fancyfoot[C]{\small\thepage}
\renewcommand{\headrulewidth}{0.4pt}
\setlength{\headheight}{15pt}

%% --- Codigo-fonte -------------------------------------------
\usepackage{listings}
\lstset{
  basicstyle      = \ttfamily\footnotesize,
  backgroundcolor = \color{almgray},
  frame           = single,
  framerule       = 0.4pt,
  rulecolor       = \color{almblue!30},
  numbers         = left,
  numberstyle     = \tiny\color{almgraytext},
  breaklines      = true,
  language        = Python,
  keywordstyle    = \color{almblue}\bfseries,
  commentstyle    = \color{almgreen}\itshape,
  stringstyle     = \color{almred},
  showstringspaces= false,
  captionpos      = b,
}

%% --- Caixas de destaque -------------------------------------
\usepackage{tcolorbox}
\tcbuselibrary{skins, breakable}
\newtcolorbox{infobox}[1][]{
  colback=almgray, colframe=almblue!60,
  fonttitle=\bfseries\small, title=#1, breakable,
  left=6pt, right=6pt, top=6pt, bottom=6pt,
}
\newtcolorbox{resultbox}[1][]{
  colback=almgreen!5, colframe=almgreen!60,
  fonttitle=\bfseries\small, title=#1, breakable,
}
\newtcolorbox{warnbox}[1][]{
  colback=almred!5, colframe=almred!50,
  fonttitle=\bfseries\small, title=#1, breakable,
}

%% --- Macros de conveniência ---------------------------------
\newcommand{\R}{\mathbb{R}}
\newcommand{\E}{\mathbb{E}}
\newcommand{\brlM}[1]{R\$\,#1\,M}
\newcommand{\brlBi}[1]{R\$\,#1\,bi}

%% --- Caminho das figuras ------------------------------------
%%     Compilar em docs/; charts estão em ../outputs/charts/
\graphicspath{{../outputs/charts/}}

%% =============================================================
\begin{document}
%% =============================================================

%% ---- CAPA --------------------------------------------------
\begin{titlepage}
  \centering
  \vspace*{1.5cm}

  {\large\bfseries\color{almblue}
    Universidade de São Paulo --- Escola Politécnica\\[0.4em]
    Programa de Pós-Graduação em Engenharia Elétrica (PPGEE)\\[0.3em]
    Área de Concentração: Engenharia de Sistemas
  }

  \vspace{2.5cm}

  \rule{\textwidth}{1.5pt}\\[0.6em]
  {\LARGE\bfseries
    Otimização Estocástica Bancária do NIM\\[0.4em]
    para \textit{Guidance} ALM-ALCO\\[0.4em]
    sob Contornos Regulatórios\\[0.4em]
    em Mercados Completos
  }\\[0.6em]
  {\large\itshape
    Relatório Técnico de Implementação --- \textit{Baseline} Quantitativo
  }\\[0.4em]
  \rule{\textwidth}{1.5pt}

  \vspace{2.0cm}

  \begin{tabular}{rl}
    \textbf{Candidato:}          & Walter Correa Neto \\[0.4em]
    \textbf{Orientador:}         & Prof.\ Dr.\ Oswaldo Luiz do Valle Costa \\[0.4em]
    \textbf{Departamento:}       & Eng.\ de Telecomunicações e Controle (EPUSP) \\[0.4em]
    \textbf{Tema:}               & Otimização Estocástica \\[0.4em]
    \textbf{Instituição-modelo:} & Banco Modelo S.A.\ (fictício) \\[0.4em]
    \textbf{Data-base:}          & Dezembro de 2024 \\[0.4em]
    \textbf{Versão:}             & 1.0 --- \textit{Baseline} QP \\
  \end{tabular}

  \vfill
  {\normalsize São Paulo, \the\year}
\end{titlepage}

%% ---- RESUMO ------------------------------------------------
\begin{abstract}
\noindent
Este relatório documenta a implementação de um sistema computacional de
\textit{Asset \& Liability Management} (ALM) bancário desenvolvido como
\textit{baseline} de pesquisa para o doutorado em Engenharia de Sistemas
(PPGEE-Poli/USP), sob orientação do Prof.~Dr.\ Oswaldo Luiz do Valle Costa.
O sistema integra oito módulos: modelagem do balanço patrimonial, geração de
cenários de taxa de juros via Nelson-Siegel \citep{NelsonSiegel1987}, motor de
fluxos de caixa, otimizador quadrático (QP) de alocação de carteira, gestão
de liquidez regulatória (LCR e NSFR) \citep{BCBS238,BCBS295} e análise de
capital econômico via RAROC \citep{BCBS2006}.
Para o Banco Modelo S.A.\ (ativo total de \brlBi{532}), o sistema produz um
NII anual de \brlM{61.026} no cenário-base, identifica um \textit{duration
gap} de 1,742~anos em posição \textit{long}, e demonstra que a fronteira
eficiente NII $\times$ risco de taxa é praticamente plana --- o custo de
reduzir o DV01 \textit{gap} de \brlM{92,7}/bp para \brlM{44}/bp representa
apenas 0,54\% do NII total.
A análise de RAROC revela que quatro segmentos operam abaixo do
\textit{hurdle rate} de 15\%, evidenciando subsídio implícito de crédito
direcionado \citep{Bonomo2016}.
O \textit{baseline} constitui referência quantitativa para a proposta de
modelo de otimização estocástica multi-período, objeto central da tese.

\medskip
\noindent\textbf{Palavras-chave:} ALM bancário; NIM; ALCO; otimização
estocástica; IRRBB; LCR; NSFR; RAROC; capital econômico; Nelson-Siegel;
programação quadrática; mercados completos.
\end{abstract}

\newpage
\tableofcontents
\newpage
\listoffigures
\newpage
\listoftables
\newpage

%% =============================================================
\section{Introdução}
\label{sec:intro}
%% =============================================================

A gestão de ativos e passivos (\textit{Asset \& Liability Management} --- ALM)
constitui uma das funções centrais da administração bancária moderna. O Comitê
de Ativos e Passivos (ALCO --- \textit{Asset \& Liability Committee}) é o órgão
colegiado responsável por traduzir as decisões de ALM em \textit{guidance}
operacional: limites de \textit{duration}, metas de NIM
(\textit{Net Interest Margin}), estratégias de hedge e alocação de capital.

Do ponto de vista quantitativo, este problema é intrinsecamente multi-objetivo
e multi-período: o banco deseja simultaneamente maximizar o NIM, controlar o
risco de taxa (IRRBB), manter \textit{buffers} de liquidez (LCR e NSFR) e
preservar adequação de capital \citep{BCBS368}. Estas dimensões são
frequentemente conflitantes, criando um espaço de decisão rico para
investigação científica no campo da Engenharia de Sistemas.

\subsection{Motivação e Contexto Acadêmico}

A literatura sobre ALM bancário apresenta três gerações de modelos:

\begin{enumerate}
  \item \textbf{Modelos determinísticos estáticos} \citep{Klein1971,Monti1972}:
    formulação de lucro máximo com \textit{spread} ótimo, ignorando dinâmica
    temporal e incerteza de taxas.
  \item \textbf{Modelos de programação estocástica}
    \citep{KusyZiemba1986,KouwenbergZenios2001}: incorporam cenários de taxa
    e restrições de liquidez em horizonte multi-período.
  \item \textbf{Modelos de otimização robusta e controle estocástico}
    \citep{BirgeLouveaux2011,Fabozzi2007}: tratam explicitamente a
    ambiguidade distribucional dos cenários de mercado.
\end{enumerate}

A prática bancária brasileira ainda opera predominantemente com a primeira
geração, usando \textit{gap reports} e análise de sensibilidade de forma
não integrada. Esta pesquisa visa construir um \textit{baseline} da segunda
geração e propor extensões para a terceira, com foco no problema de
otimização do NIM sob contornos regulatórios em mercados completos ---
tema alinhado à área de concentração em Engenharia de Sistemas do
PPGEE-Poli/USP e à experiência do orientador em controle estocástico
\citep{Costa2005}.

\subsection{Objetivos}

\begin{enumerate}
  \item Implementar um sistema computacional completo de ALM para um banco de
    varejo de médio-grande porte com calibração realista para o mercado
    brasileiro (dezembro de 2024).
  \item Estabelecer um \textit{baseline} metodológico robusto baseado em
    Programação Quadrática (QP), documentado e reprodutível.
  \item Identificar as limitações estruturais do \textit{baseline} que motivem
    a contribuição científica original da tese.
\end{enumerate}

\subsection{Estrutura do Documento}

O relatório está organizado em oito seções: modelagem do balanço
(\S\ref{sec:balanco}), cenários de mercado (\S\ref{sec:cenarios}), motor de
fluxos de caixa (\S\ref{sec:cashflow}), otimizador QP (\S\ref{sec:optimizer}),
módulos de liquidez e capital (\S\ref{sec:liqcap}), perspectivas de pesquisa
(\S\ref{sec:perspectivas}), arquitetura (\S\ref{sec:arquitetura}) e conclusões
(\S\ref{sec:conclusoes}).

%% =============================================================
\section{Modelagem do Balanço Patrimonial}
\label{sec:balanco}
%% =============================================================

\subsection{Estrutura Geral}

O Banco Modelo S.A.\ é uma instituição financeira fictícia calibrada para
representar um banco de varejo brasileiro de médio-grande porte, com
data-base de 31 de dezembro de 2024.
A Tabela~\ref{tab:balanco_resumo} apresenta os agregados patrimoniais.

\begin{table}[H]
\centering
\caption{Resumo do Balanço Patrimonial --- Banco Modelo S.A.\ (R\$ milhões)}
\label{tab:balanco_resumo}
\begin{tabular}{lrr}
\toprule
\textbf{Item} & \textbf{R\$ M} & \textbf{\% Ativo} \\
\midrule
\multicolumn{3}{l}{\textit{Ativo}} \\
\quad Crédito Total                  & 342.800 & 64,44\% \\
\quad Títulos e Valores Mobiliários  & 114.600 & 21,54\% \\
\quad Disponibilidades e Compulsório &  40.500 &  7,61\% \\
\quad Operações Interbancárias       &  11.200 &  2,11\% \\
\quad Outros Ativos                  &  22.900 &  4,30\% \\
\midrule
\textbf{Total Ativo}                 & \textbf{532.000} & \textbf{100,00\%} \\
\midrule
\multicolumn{3}{l}{\textit{Passivo}} \\
\quad CDB (várias séries)            & 116.300 & 21,86\% \\
\quad Depósitos à Vista e Poupança   &  88.500 & 16,63\% \\
\quad Letras Financeiras             &  50.000 &  9,40\% \\
\quad LCI / LCA                      &  40.200 &  7,56\% \\
\quad \textit{Funding} Externo       &  20.800 &  3,91\% \\
\quad Outros Passivos                &  23.700 &  4,46\% \\
\midrule
\textbf{Total Passivo}               & \textbf{339.500} & \textbf{63,82\%} \\
\midrule
\textbf{Patrimônio Líquido}          & \textbf{192.100} & \textbf{36,11\%} \\
\bottomrule
\end{tabular}
\fonte{Elaborado pelo autor.}
\end{table}

\begin{figure}[H]
  \centering
  \includegraphics[width=\textwidth]{01_composicao_balanco.png}
  \caption{Composição do balanço patrimonial por categoria de ativo e passivo.
           O crédito representa 64,44\% do ativo total; CDB e depósitos
           respondem pela maior fatia do \textit{funding}.}
  \label{fig:composicao_balanco}
  \fonte{Elaborado pelo autor.}
\end{figure}

\subsection{Indicadores de Capital Regulatório}

O sistema calcula automaticamente os indicadores de adequação de capital
conforme a Resolução BCB 4.958/2021 e Basileia III \citep{BCBS2006}.
Os resultados estão na Tabela~\ref{tab:capital_regulatorio}.

\begin{table}[H]
\centering
\caption{Indicadores de Capital Regulatório --- Data-base Dez/2024}
\label{tab:capital_regulatorio}
\begin{tabular}{lrrl}
\toprule
\textbf{Indicador} & \textbf{Valor} & \textbf{Mínimo} & \textbf{Status} \\
\midrule
RWA Total                     & R\$ 265.690 M & ---    & --- \\
Patrimônio de Referência (PR) & R\$ 37.800 M  & ---    & --- \\
Índice de Basileia (IB)       & 14,23\%       & 10,5\% & \textcolor{almgreen}{\textbf{OK}} \\
CET1 \textit{Ratio}           & 10,73\%       & 7,0\%  & \textcolor{almgreen}{\textbf{OK}} \\
\textit{Leverage}             & 2,77$\times$  & ---    & --- \\
Crédito / Ativo               & 64,44\%       & ---    & --- \\
\bottomrule
\end{tabular}
\fonte{Elaborado pelo autor.}
\end{table}

\begin{figure}[H]
  \centering
  \includegraphics[width=\textwidth]{02_rwa_capital.png}
  \caption{Distribuição do RWA por categoria de ativo e posição de capital.
           O crédito a pessoas físicas e jurídicas representa a maior parcela
           do ativo ponderado pelo risco.}
  \label{fig:rwa_capital}
  \fonte{Elaborado pelo autor.}
\end{figure}

\subsection{Mix de Indexadores}

A estrutura de indexação do balanço determina diretamente a sensibilidade
do NIM a diferentes cenários de taxa.
A Figura~\ref{fig:mix_indexadores} apresenta o mix ativo vs.\ passivo.

\begin{figure}[H]
  \centering
  \includegraphics[width=\textwidth]{03_mix_indexadores.png}
  \caption{Mix de indexadores --- ativo vs.\ passivo. No ativo, 58,2\% é
           indexado ao CDI; no passivo, 60,5\%. A posição em IPCA é
           assimétrica: o passivo detém proporção maior, criando exposição
           líquida a inflação.}
  \label{fig:mix_indexadores}
  \fonte{Elaborado pelo autor.}
\end{figure}

%% =============================================================
\section{Cenários de Mercado}
\label{sec:cenarios}
%% =============================================================

\subsection{Modelagem da Curva de Juros: Nelson-Siegel}

As curvas de juros são geradas via modelo de \citet{NelsonSiegel1987},
calibrado com dados de mercado de dezembro de 2024:

\begin{equation}
  r(\tau) = \beta_0
    + \beta_1 \cdot \frac{1 - e^{-\tau/\lambda}}{\tau/\lambda}
    + \beta_2 \cdot \left[
        \frac{1 - e^{-\tau/\lambda}}{\tau/\lambda} - e^{-\tau/\lambda}
      \right]
  \label{eq:nelson_siegel}
\end{equation}

\noindent onde $r(\tau)$ é a taxa \textit{spot} ao prazo $\tau$ (anos),
$\beta_0$ é o fator de nível (taxa de longo prazo), $\beta_1$ o fator de
inclinação, $\beta_2$ o fator de curvatura e $\lambda$ o parâmetro de
decaimento. O modelo foi proposto como forma parcimoniosa de representar
toda a estrutura a termo com apenas quatro parâmetros
\citep{NelsonSiegel1987}.

Os parâmetros calibrados para o cenário-base estão na
Tabela~\ref{tab:nelson_siegel}.

\begin{table}[H]
\centering
\caption{Parâmetros Nelson-Siegel --- Calibração Dez/2024 (Cenário-Base)}
\label{tab:nelson_siegel}
\begin{tabular}{lrl}
\toprule
\textbf{Parâmetro} & \textbf{Valor} & \textbf{Interpretação} \\
\midrule
$\beta_0$ & 12,80\% & Nível de longo prazo \\
$\beta_1$ & $-1,30\%$ & Inclinação (spread curto $-$ longo) \\
$\beta_2$ & $-2,50\%$ & Curvatura \\
$\lambda$ & 3,50 & Parâmetro de decaimento \\
\midrule
SELIC (âncora curta) & 11,75\% & Taxa \textit{overnight} \\
\bottomrule
\end{tabular}
\fonte{Elaborado pelo autor.}
\end{table}

\subsection{Cenários de Estresse}

São gerados três cenários de taxa de juros:

\begin{description}
  \item[\textbf{Base}] Curva calibrada a Dez/2024, SELIC = 11,75\% a.a.
  \item[\textbf{Stress Alta}] Choque paralelo de $+200$ bps na curva DI.
  \item[\textbf{Stress Baixa}] Choque paralelo de $-250$ bps na curva DI.
\end{description}

Os cenários seguem a metodologia BCBS~368 \citep{BCBS368} para IRRBB.

\begin{figure}[H]
  \centering
  \includegraphics[width=\textwidth]{04_curvas_juros_cenarios.png}
  \caption{Curvas de juros DI e IPCA para os três cenários. A curva base
           apresenta leve inversão para prazos intermediários, refletindo
           expectativa de queda gradual da SELIC.}
  \label{fig:curvas_juros}
  \fonte{Elaborado pelo autor.}
\end{figure}

\subsection{Spreads de Crédito}

Os \textit{spreads} de crédito por segmento são calibrados com referências
de mercado e sensibilizados nos cenários de estresse.

\begin{figure}[H]
  \centering
  \includegraphics[width=\textwidth]{05_spreads_credito.png}
  \caption{\textit{Spreads} de crédito por segmento nos três cenários
           (escala logarítmica). No cenário de estresse de alta, os
           \textit{spreads} ampliam-se, refletindo deterioração das
           condições de crédito e aumento do prêmio de risco.}
  \label{fig:spreads_credito}
  \fonte{Elaborado pelo autor.}
\end{figure}

%% =============================================================
\section{Motor de Fluxos de Caixa e NIM}
\label{sec:cashflow}
%% =============================================================

\subsection{Gap de Reprecificação}

O \textit{gap} de reprecificação mensura a diferença entre ativos e passivos
que repreci\-ficam dentro de cada \textit{bucket} temporal, conforme
\citet{SaundersWalter2011} e as diretrizes do BCBS~368 \citep{BCBS368}.
Os resultados estão na Tabela~\ref{tab:gap_reprecificacao} e na
Figura~\ref{fig:gap_reprecificacao}.

\begin{table}[H]
\centering
\caption{Gap de Reprecificação por \textit{Bucket} Temporal (R\$ milhões)}
\label{tab:gap_reprecificacao}
\begin{tabular}{lrrrrr}
\toprule
\textbf{Bucket} & \textbf{Ativo} & \textbf{Passivo} & \textbf{Gap} &
  \textbf{Gap Cum.} & \textbf{DV01 Gap} \\
\midrule
O/N--1M  & 200.800 & 246.000 & $-$45.200 & $-$45.200 & $+$0,01 \\
1M--3M   &       0 &       0 &         0 & $-$45.200 & $+$0,00 \\
3M--6M   &       0 &       0 &         0 & $-$45.200 & $+$0,00 \\
6M--1A   &  34.000 &  23.000 & $+$11.000 & $-$34.200 & $+$0,20 \\
1A--2A   &  40.900 &   9.800 & $+$31.100 &  $-$3.100 & $+$4,20 \\
2A--3A   &  40.800 &  16.200 & $+$24.600 & $+$21.500 & $+$5,74 \\
3A--5A   & 102.500 &  12.500 & $+$90.000 & $+$111.500& $+$29,73 \\
5A--7A   &  28.000 &  32.000 &  $-$4.000 & $+$107.500& $-$0,65 \\
7A--10A  &       0 &       0 &         0 & $+$107.500& $+$0,00 \\
10A+     &  85.000 &       0 & $+$85.000 & $+$192.500& $+$53,44 \\
\midrule
\textbf{Total} & \textbf{532.000} & \textbf{339.500} & & &
  \textbf{$+$92,67} \\
\bottomrule
\end{tabular}
\nota{DV01 Gap em R\$\,M/bp. Gap positivo indica posição \textit{long} no
      \textit{bucket}.}
\fonte{Elaborado pelo autor.}
\end{table}

\begin{infobox}[Interpretação do Gap de Reprecificação]
O gap negativo em O/N--1M ($-$R\$\,45,2\,bi) é típico de banco de varejo:
os passivos CDI (depósitos, CDB de curto prazo) repreci\-ficam antes dos
ativos de crédito de médio prazo. O gap cumulativo inverte no bucket 2A--3A
e atinge R\$\,192,5\,bi positivo em 10A+, sinalizando posição \textit{long}
de \textit{duration}.
\end{infobox}

\begin{figure}[H]
  \centering
  \includegraphics[width=\textwidth]{06_gap_reprecificacao.png}
  \caption{Gap de reprecificação por \textit{bucket}: barras representam
           volumes de ativo (azul) e passivo (vermelho) por prazo; a linha
           tracejada representa o gap cumulativo.}
  \label{fig:gap_reprecificacao}
  \fonte{Elaborado pelo autor.}
\end{figure}

\subsection{Projeção de NII e NIM}

O NII (\textit{Net Interest Income}) é calculado mensalmente para cada produto.
A Tabela~\ref{tab:nii_resumo} apresenta os resultados anuais por cenário.

\begin{table}[H]
\centering
\caption{Projeção de NII --- 12 meses, por cenário (R\$ milhões)}
\label{tab:nii_resumo}
\begin{tabular}{lrrrrr}
\toprule
\textbf{Cenário} & \textbf{NII Acum.} & \textbf{NIM Médio} &
  \textbf{Receita} & \textbf{Custo} & \textbf{$\Delta$ vs.\ Base} \\
\midrule
Base          & 61.026 & 11,47\% &  95.703 & 34.677 & --- \\
Stress Alta   & 60.668 & 11,40\% & 104.595 & 43.927 & $-$358 \\
Stress Baixa  & 61.353 & 11,53\% &  86.970 & 25.617 & $+$327 \\
\bottomrule
\end{tabular}
\fonte{Elaborado pelo autor.}
\end{table}

\begin{figure}[H]
  \centering
  \includegraphics[width=\textwidth]{07_nii_cenarios.png}
  \caption{NII mensal acumulado para os três cenários de taxa de juros.
           A estabilidade relativa reflete o alto grau de indexação ao CDI
           em ambos os lados do balanço.}
  \label{fig:nii_cenarios}
  \fonte{Elaborado pelo autor.}
\end{figure}

\begin{figure}[H]
  \centering
  \includegraphics[width=\textwidth]{08_nim_cenarios.png}
  \caption{NIM (\% a.a.) por cenário ao longo dos 12 meses de projeção.
           O NIM mantém-se entre 11,40\% e 11,53\% entre os cenários de
           alta e baixa, respectivamente.}
  \label{fig:nim_cenarios}
  \fonte{Elaborado pelo autor.}
\end{figure}

\subsection{Duration e DV01}

A \textit{duration} modificada é calculada analiticamente por produto
\citep{Fabozzi1996}:

\begin{equation}
  D_{\text{mod}} = \frac{D_{\text{Mac}}}{1 + y/m}
  \label{eq:duration_mod}
\end{equation}

\noindent onde $D_{\text{Mac}}$ é a \textit{duration} de Macaulay,
$y$ a \textit{yield} e $m$ a frequência de cupom.
O DV01 (\textit{Dollar Value of 1 basis point}) é:

\begin{equation}
  \text{DV01} = \text{Saldo} \times D_{\text{mod}} \times 0{,}0001
  \label{eq:dv01}
\end{equation}

Os resultados consolidados são:
\begin{itemize}
  \item \textit{Duration} ativo: 2,047 anos;
  \item \textit{Duration} passivo: 0,478 anos;
  \item \textbf{\textit{Duration gap}: 1,742 anos} (posição \textit{long});
  \item DV01 ativo: R\$\,108,9\,M/bp;
  \item DV01 passivo: R\$\,16,2\,M/bp;
  \item \textbf{DV01 gap: R\$\,92,7\,M/bp}.
\end{itemize}

\begin{figure}[H]
  \centering
  \includegraphics[width=\textwidth]{09_duration_dv01.png}
  \caption{DV01 por produto e \textit{duration} modificada média. A carteira
           de crédito imobiliário e os títulos NTN-B longa concentram o maior
           DV01 do ativo.}
  \label{fig:duration_dv01}
  \fonte{Elaborado pelo autor.}
\end{figure}

\begin{figure}[H]
  \centering
  \includegraphics[width=\textwidth]{10_sensibilidade_nii.png}
  \caption{Sensibilidade do NII a choques paralelos de $-300$ a $+300$ bps,
           metodologia IRRBB \citep{BCBS368}. A linearidade próxima reflete
           a predominância de instrumentos CDI, que repreci\-ficam rapidamente.}
  \label{fig:sensibilidade_nii}
  \fonte{Elaborado pelo autor.}
\end{figure}

\begin{warnbox}[Risco de Taxa de Juros --- Posição LONG]
Com \textit{duration gap} de 1,742 anos, o banco perde valor econômico
quando taxas sobem. Um choque de $+100$ bps reduz o NII em aproximadamente
R\$\,180\,M (sensibilidade linear). A posição está dentro do limite de
$\pm 2{,}0$ anos (IRRBB), mas próxima do teto, justificando hedge ativo.
\end{warnbox}

%% =============================================================
\section{Otimizador ALM: \textit{Baseline} QP}
\label{sec:optimizer}
%% =============================================================

\subsection{Formulação do Problema}

O otimizador ALM resolve o seguinte problema de Programação Quadrática (QP)
\citep{Nocedal2006}:

\begin{equation}
\begin{aligned}
  \max_{\mathbf{x}} \quad
    & \text{NII}(\mathbf{x})
      - \lambda \cdot \bigl[\text{DV01\_gap}(\mathbf{x})\bigr]^2 \\[6pt]
  \text{s.a.} \quad
    & \text{LCR}(\mathbf{x}) \geq 1{,}00 \\
    & \text{NSFR}(\mathbf{x}) \geq 1{,}00 \\
    & \text{IB}(\mathbf{x}) \geq 10{,}5\% \\
    & \bigl|\text{DurGap}(\mathbf{x})\bigr| \leq 2{,}0 \text{ anos} \\
    & \sum_s x_{\text{cred},s} + \sum_t \max(0,\,x_{\text{tvm},t})
        = \sum_f x_{\text{cap},f} \\
    & 0 \leq x_{\text{cred},s} \leq D_s(\text{spread}_s) \\
    & 0 \leq x_{\text{cap},f} \leq V_f^{\max} \\
    & x_{\text{tvm},t} \in \bigl[-S_t,\,V_t^{\text{disp}}\bigr] \\
    & x_{\text{hedge},h} \in \bigl[-N_h^{\max},\,N_h^{\max}\bigr]
\end{aligned}
\label{eq:qp}
\end{equation}

\noindent onde $\lambda \geq 0$ é o parâmetro de aversão ao risco de taxa,
$D_s(\cdot)$ é a curva de demanda do segmento $s$, $V_f^{\max}$ o volume
máximo de captação e $N_h^{\max}$ o nocional máximo do instrumento de hedge $h$.

\subsection{Variáveis de Decisão}

O vetor $\mathbf{x} \in \R^{28}$, com $N = 7 + 6 + 8 + 7 = 28$ componentes:

\begin{equation}
  \mathbf{x} = \Bigl[
    \underbrace{x_{\text{cred},1},\ldots,x_{\text{cred},7}}_{N_c=7},\;
    \underbrace{x_{\text{cap},1},\ldots,x_{\text{cap},6}}_{N_f=6},\;
    \underbrace{x_{\text{tvm},1},\ldots,x_{\text{tvm},8}}_{N_t=8},\;
    \underbrace{x_{\text{hdg},1},\ldots,x_{\text{hdg},7}}_{N_h=7}
  \Bigr]^{\!\top}
  \label{eq:vetor_x}
\end{equation}

\subsection{Modelagem da Demanda de Crédito}

A curva de demanda de crédito por segmento é modelada como:

\begin{equation}
  V_s(\text{spread}) = V_s^0 \cdot
    \exp\!\bigl(-\varepsilon_s \cdot (\text{spread} - \text{spread}_s^0)\bigr)
  \label{eq:demanda_credito}
\end{equation}

\noindent onde $\varepsilon_s > 0$ é a elasticidade \textit{spread}-volume.
Esta especificação é consistente com o modelo Klein-Monti
\citep{Klein1971,Monti1972} e com evidências empíricas de
\citet{Coelho2012} para o Brasil.

\subsection{Modelagem do Custo de Captação}

\begin{equation}
  c_f(v) = c_f^0 + \gamma_f \cdot
    \left(\frac{v}{v_f^{\text{ref}}}\right)^{\!\delta_f}
  \label{eq:custo_captacao}
\end{equation}

\noindent com $\delta_f = 0{,}5$: custo cresce com a raiz quadrada do volume.

\subsection{Solver e Implementação}

O problema \eqref{eq:qp} é resolvido via \texttt{scipy.optimize.minimize}
com método SLSQP (\textit{Sequential Least Squares Programming})
\citep{Kraft1988}.

\begin{lstlisting}[caption={Estrutura central do otimizador
                   (\texttt{ALMOptimizer.otimizar})},
                   label={lst:optimizer}]
resultado = minimize(
    fun    = lambda x: -NII(x) + lambda_risco * DV01_gap(x)**2,
    x0     = 0.2 * upper_bounds,
    method = "SLSQP",
    bounds = Bounds(lb, ub),
    constraints = [
        {"type":"ineq", "fun": lambda x: LCR(x)    - 1.00 },
        {"type":"ineq", "fun": lambda x: NSFR(x)   - 1.00 },
        {"type":"ineq", "fun": lambda x: IB(x)     - 0.105},
        {"type":"ineq", "fun": lambda x: 2.0 - abs(DurGap(x))},
        {"type":"eq",   "fun": lambda x: sum(usos) - sum(fontes)},
    ],
    options={"maxiter": 1000, "ftol": 1e-8},
)
\end{lstlisting}

\subsection{Fronteira Eficiente}

A fronteira eficiente é traçada variando $\lambda \in [10^{-4}, 1{,}0]$
com 12 pontos em escala logarítmica. Os resultados estão na
Tabela~\ref{tab:fronteira_eficiente} e na Figura~\ref{fig:fronteira_eficiente}.

\begin{table}[H]
\centering
\caption{Fronteira Eficiente --- NII $\times$ Risco de Taxa}
\label{tab:fronteira_eficiente}
\begin{tabular}{rrrrrr}
\toprule
$\lambda$ & \textbf{NII Total (R\$M)} & \textbf{DV01 Gap} &
  \textbf{Dur.\ Gap (a)} & \textbf{LCR} & \textbf{IB} \\
\midrule
0,0001 & 61.830 & 90,4 & 1,58 & 359\% & 13,8\% \\
0,0012 & 61.826 & 88,0 & 1,54 & 359\% & 13,8\% \\
0,0066 & 61.823 & 75,5 & 1,32 & 359\% & 13,8\% \\
0,0152 & 61.802 & 55,6 & 0,97 & 359\% & 13,8\% \\
0,0811 & 61.787 & 48,0 & 0,84 & 360\% & 13,8\% \\
1,0000 & 61.499 & 43,9 & 0,77 & 362\% & 13,9\% \\
\bottomrule
\end{tabular}
\fonte{Elaborado pelo autor.}
\end{table}

\begin{figure}[H]
  \centering
  \includegraphics[width=\textwidth]{11_fronteira_eficiente.png}
  \caption{Fronteira eficiente NII $\times$ DV01 gap. Os três pontos
           destacados correspondem às estratégias de máximo NII
           ($\lambda = 0{,}0001$), equilíbrio ($\lambda = 0{,}01$)
           e mínimo risco ($\lambda = 1{,}0$).}
  \label{fig:fronteira_eficiente}
  \fonte{Elaborado pelo autor.}
\end{figure}

\begin{resultbox}[Achado Central --- Fronteira Eficiente Quase Plana]
Reduzir o DV01 gap de \brlM{90,4}/bp para \brlM{43,9}/bp custa apenas
\brlM{331} de NII, ou 0,54\% do NII total. O custo do hedge de taxa é
relativamente baixo. A persistência de bancos brasileiros com
\textit{duration gap} elevado sugere restrições institucionais, contábeis
(EVE vs.\ NIM) ou de incentivos gerenciais não capturadas neste modelo
--- motivação direta para a tese.
\end{resultbox}

\subsection{Alocação Ótima de Crédito e Captação}

\begin{figure}[H]
  \centering
  \includegraphics[width=\textwidth]{12_alocacao_credito.png}
  \caption{Volume de crédito novo por segmento para as três estratégias.
           O otimizador aloca sistematicamente em crédito pessoal, veículos
           e PJ, nunca em imobiliário ou rural (RAROC negativo no modelo).}
  \label{fig:alocacao_credito}
  \fonte{Elaborado pelo autor.}
\end{figure}

\begin{figure}[H]
  \centering
  \includegraphics[width=\textwidth]{13_mix_captacao.png}
  \caption{Mix de captação nova ótima por instrumento. A preferência por
           CDB CDI longo e LF IPCA reflete menor custo e contribuição
           positiva para o NSFR.}
  \label{fig:mix_captacao}
  \fonte{Elaborado pelo autor.}
\end{figure}

\begin{figure}[H]
  \centering
  \includegraphics[width=\textwidth]{14_restricoes_regulatorias.png}
  \caption{Atendimento às restrições regulatórias pós-otimização para as
           três estratégias. O NSFR é a restrição ativa (\textit{binding})
           em todas as soluções; LCR e Índice de Basileia mantêm folga
           substancial.}
  \label{fig:restricoes}
  \fonte{Elaborado pelo autor.}
\end{figure}

\begin{figure}[H]
  \centering
  \includegraphics[width=\textwidth]{15_hedge_dv01.png}
  \caption{Posição nocional de hedge e DV01 gap resultante por estratégia.
           A estratégia de mínimo risco utiliza o limite máximo de nocional
           em futuros DI e \textit{swaps} DI$\times$Pré.}
  \label{fig:hedge_dv01}
  \fonte{Elaborado pelo autor.}
\end{figure}

%% =============================================================
\section{Módulos de Liquidez e Capital}
\label{sec:liqcap}
%% =============================================================

\subsection{Stress Test de Liquidez --- LCR}

O LCR (\textit{Liquidity Coverage Ratio}) é calculado conforme a Circular BCB
3.644/2013 e Basileia III \citep{BCBS238}:

\begin{equation}
  \text{LCR} = \frac{\text{HQLA}}{\text{NCO}_{30}} \geq 1{,}00, \quad
  \text{NCO} = \text{Out}_{30d} -
    \min\!\bigl(\text{In}_{30d},\; 0{,}75\cdot\text{Out}_{30d}\bigr)
  \label{eq:lcr}
\end{equation}

\begin{table}[H]
\centering
\caption{Resultado do Stress Test LCR --- 30 dias}
\label{tab:lcr_stress}
\begin{tabular}{lrr}
\toprule
\textbf{Componente} & \textbf{R\$ M} & \textbf{Detalhe} \\
\midrule
HQLA Nível 1 (0\% \textit{haircut})  & 152.200 & LFT, Reservas BCB \\
HQLA Nível 2A ($-$15\%)              &   9.350 & CDB outros bancos \\
HQLA Nível 2B ($-$25\%)              &   2.325 & Debêntures A \\
\midrule
\textbf{HQLA Total}                  & \textbf{163.875} & \\
Outflows brutos (30d)                &  22.931 & Estresse Basileia III \\
Inflows (cap 75\%)                   &  10.427 & \\
\textbf{NCO}                         & \textbf{12.504} & \\
\midrule
\textbf{LCR = HQLA / NCO}           & \textbf{1.311\%} &
  \textcolor{almgreen}{$\checkmark$ Aprovado} \\
\bottomrule
\end{tabular}
\fonte{Elaborado pelo autor.}
\end{table}

\begin{figure}[H]
  \centering
  \includegraphics[width=\textwidth]{16_stress_test_lcr.png}
  \caption{Decomposição do \textit{stress test} LCR: composição HQLA por
           nível regulatório e principais \textit{outflows} líquidos em
           cenário de estresse. LCR de 1.311\% indica buffer muito acima
           do mínimo exigido.}
  \label{fig:stress_lcr}
  \fonte{Elaborado pelo autor.}
\end{figure}

\subsection{NSFR e Projeção Dinâmica}

O NSFR (\textit{Net Stable Funding Ratio}) \citep{BCBS295} avalia a
sustentabilidade estrutural do \textit{funding}:

\begin{equation}
  \text{NSFR} = \frac{\text{ASF}}{\text{RSF}} \geq 1{,}00
  \label{eq:nsfr}
\end{equation}

\noindent O resultado base é NSFR = 160,9\%
(ASF = R\$\,442,6\,bi; RSF = R\$\,275,0\,bi).

\begin{figure}[H]
  \centering
  \includegraphics[width=\textwidth]{17_projecao_lcr_nsfr.png}
  \caption{Projeção dinâmica de LCR e NSFR ao longo de 12 meses para os
           três cenários de crescimento. O NSFR comprime-se mais rapidamente
           e torna-se a restrição ativa (\textit{binding}) no otimizador.}
  \label{fig:projecao_lcr_nsfr}
  \fonte{Elaborado pelo autor.}
\end{figure}

\subsection{Custo do Buffer de Liquidez}

O custo de oportunidade de cada produto HQLA é:

\begin{equation}
  \text{Custo}_{\text{HQLA}} =
    \text{Saldo} \times \bigl(r_{\text{alt}} - r_{\text{HQLA}}\bigr)
  \label{eq:custo_hqla}
\end{equation}

\noindent onde $r_{\text{alt}}$ é a taxa da melhor aplicação alternativa
(capital de giro PJ ao CDI + 4,5\%). O custo total é \brlM{9.400}/ano.

\begin{figure}[H]
  \centering
  \includegraphics[width=\textwidth]{18_custo_buffer_hqla.png}
  \caption{Custo de carregamento por produto HQLA (R\$\,M/ano). LFT e
           compulsório respondem pela maior parte do custo de oportunidade
           total de \brlM{9.400}/ano.}
  \label{fig:custo_hqla}
  \fonte{Elaborado pelo autor.}
\end{figure}

\subsection{Capital Econômico e RAROC}

O capital econômico (CE) é calculado via fórmula IRB \citep{BCBS2006}:

\begin{equation}
  \text{CE}_s = \text{EAD}_s \cdot \text{LGD}_s \cdot
    \left[
      \Phi\!\left(
        \frac{\Phi^{-1}(\text{PD}_s) + \sqrt{\rho_s}\;\Phi^{-1}(\alpha)}
             {\sqrt{1-\rho_s}}
      \right) - \text{PD}_s
    \right] \cdot \text{MA}
  \label{eq:ce_irb}
\end{equation}

\noindent onde $\Phi$ é a CDF normal padrão, $\alpha = 99{,}9\%$ o nível de
confiança, $\rho_s$ a correlação sistêmica e MA o ajuste de maturidade.
O RAROC é calculado conforme \citet{MertonPerold1993}:

\begin{equation}
  \text{RAROC}_s =
    \frac{\text{Receita}_s - \text{Custo Funding}_s
          - \text{PE}_s - \text{Custo Op}_s}
         {\text{CE}_s}
  \label{eq:raroc}
\end{equation}

\begin{table}[H]
\centering
\caption{RAROC por Linha de Negócio (\textit{Hurdle Rate}: 15\%)}
\label{tab:raroc}
\small
\begin{tabular}{lrrrrrr}
\toprule
\textbf{Segmento} & \textbf{EAD} & \textbf{CE} & \textbf{NII Líq.} &
  \textbf{RAROC} & \textbf{EVA} & \textbf{OK?} \\
& (R\$\,bi) & (R\$\,bi) & (R\$\,M) & (\%) & (R\$\,M) & \\
\midrule
Cartão PF       & 19,5 & 3,70 & 32.879 &  888,1 & $+$32.324 & \textcolor{almgreen}{$\checkmark$} \\
Consignado PF   & 62,0 & 3,12 &  5.602 &  179,6 & $+$5.134  & \textcolor{almgreen}{$\checkmark$} \\
Pessoal PF      & 24,0 & 3,11 &  4.730 &  152,3 & $+$4.265  & \textcolor{almgreen}{$\checkmark$} \\
Veículos PF     & 31,0 & 2,73 &  3.802 &  139,1 & $+$3.392  & \textcolor{almgreen}{$\checkmark$} \\
Corporativo PJ  & 52,0 & 4,61 &    835 &   18,1 & $+$143    & \textcolor{almgreen}{$\checkmark$} \\
Capital Giro PJ & 38,0 & 4,31 &    397 &    9,2 & $-$249    & \textcolor{almred}{$\times$} \\
Imobiliário PF  & 85,0 & 6,16 & $-$1.424 & $-$23,1 & $-$2.348 & \textcolor{almred}{$\times$} \\
Rural           & 16,8 & 0,79 &  $-$439 & $-$55,3 &  $-$559 & \textcolor{almred}{$\times$} \\
Trade Finance   & 14,5 & 0,18 &  $-$702 & $-$392,9& $-$729  & \textcolor{almred}{$\times$} \\
\midrule
\textbf{Total}  & 342,8 & 28,71 & 45.680 & 159,1 & $+$41.373 & --- \\
\bottomrule
\end{tabular}
\fonte{Elaborado pelo autor.}
\end{table}

\begin{figure}[H]
  \centering
  \includegraphics[width=\textwidth]{19_raroc_linhas_negocio.png}
  \caption{RAROC por segmento de negócio versus \textit{hurdle rate} de 15\%.
           Segmentos à esquerda da linha vermelha destroem valor econômico
           (EVA negativo).}
  \label{fig:raroc}
  \fonte{Elaborado pelo autor.}
\end{figure}

\subsection{Pricing Mínimo via RAROC}

O \textit{spread} mínimo para atingir RAROC $\geq 15\%$ é
\citep{ResteSironi2007}:

\begin{equation}
  r_s^{\min} = r_{\text{fund}}
             + \frac{\text{PE}_s}{\text{EAD}_s}
             + c_{\text{op},s}
             + \bar{r} \cdot \frac{\text{CE}_s}{\text{EAD}_s}
  \label{eq:pricing_minimo}
\end{equation}

\begin{table}[H]
\centering
\caption{Pricing Mínimo por Segmento (RAROC $\geq$ 15\%)}
\label{tab:pricing}
\begin{tabular}{lrrrl}
\toprule
\textbf{Segmento} & \textbf{Taxa Atual} & \textbf{Taxa Mín.} &
  \textbf{Folga} & \textbf{OK?} \\
\midrule
Cartão PF       & 189,00\% & 23,24\% & $+$165,76 pp & \textcolor{almgreen}{$\checkmark$} \\
Pessoal PF      &  38,40\% & 20,63\% & $+$17,77 pp  & \textcolor{almgreen}{$\checkmark$} \\
Veículos PF     &  26,80\% & 15,86\% & $+$10,94 pp  & \textcolor{almgreen}{$\checkmark$} \\
Consignado PF   &  22,50\% & 14,22\% &  $+$8,28 pp  & \textcolor{almgreen}{$\checkmark$} \\
Corporativo PJ  &  14,55\% & 14,27\% &  $+$0,28 pp  & \textcolor{almgreen}{$\checkmark$} \\
Capital Giro PJ &  16,25\% & 16,91\% &  $-$0,66 pp  & \textcolor{almred}{$\times$} \\
Imobiliário PF  &  10,50\% & 13,26\% &  $-$2,76 pp  & \textcolor{almred}{$\times$} \\
Rural           &  10,50\% & 13,82\% &  $-$3,32 pp  & \textcolor{almred}{$\times$} \\
Trade Finance   &   7,40\% & 12,42\% &  $-$5,02 pp  & \textcolor{almred}{$\times$} \\
\bottomrule
\end{tabular}
\fonte{Elaborado pelo autor.}
\end{table}

\begin{figure}[H]
  \centering
  \includegraphics[width=\textwidth]{20_pricing_capital_economico.png}
  \caption{Comparação entre taxa atual e taxa mínima econômica por segmento,
           e capital econômico como percentual do EAD. Segmentos regulados
           (imobiliário, rural) são mantidos por exigências de crédito
           direcionado \citep{Bonomo2016}, não por retorno econômico.}
  \label{fig:pricing}
  \fonte{Elaborado pelo autor.}
\end{figure}

%% =============================================================
\section{Limitações do \textit{Baseline} e Perspectivas de Pesquisa}
\label{sec:perspectivas}
%% =============================================================

\subsection{Limitações Identificadas}

A análise do \textit{baseline} QP revela três limitações estruturais que
motivam diretamente a contribuição da tese:

\subsubsection{Fronteira Eficiente Quase Plana}

A variação de apenas 0,54\% no NIM ao longo da fronteira levanta a questão:
\textit{por que bancos brasileiros mantêm posições \textit{long} de duration
substanciais se o custo do hedge é tão baixo?} Possíveis explicações:
\begin{itemize}
  \item \textit{Hedge accounting} (CPC 48 / IFRS 9): custos de documentação
    e teste de efetividade.
  \item Assimetria EVE vs.\ NIM: gestão por NIM ignora o impacto no
    valor econômico do patrimônio \citep{BCBS368}.
  \item Fricções institucionais: CVA/DVA de derivativos, limites de nocional.
  \item Incentivos gerenciais: metas de NIM de curto prazo desincentivam
    hedges que custam margem imediata.
\end{itemize}

\subsubsection{NSFR como Restrição Dominante}

O NSFR é a restrição ativa em todas as soluções. Modelos estáticos de período
único não capturam como decisões de hoje afetam o NSFR nos próximos 12 meses
à medida que o mix de captação amadurece \citep{KoussisPadgett2016}.

\subsubsection{RAROC Negativo em Segmentos Regulados}

Imobiliário, rural e \textit{trade finance} apresentam RAROC negativo, mas
são mantidos por razões estratégicas, regulatórias (crédito direcionado) ou
de externalidades não capturadas \citep{Bonomo2016,Coelho2012}.

\subsection{Agenda de Pesquisa e Conexão com a Tese}

\begin{enumerate}
  \item \textbf{Programação Estocástica Multi-Período}
    \citep{BirgeLouveaux2011,KusyZiemba1986}: reformulação de
    \eqref{eq:qp} como problema de dois estágios:
    \begin{equation}
      \max_{\mathbf{x}_0,\,\mathbf{x}(\xi)} \;
        \E_\xi\!\bigl[\text{NIM}(\mathbf{x}(\xi))\bigr]
      \quad \text{s.a.\ restrições regulatórias } \forall\,\xi
      \label{eq:stoch_prog}
    \end{equation}
    onde $\mathbf{x}_0$ são decisões do primeiro estágio (agora) e
    $\mathbf{x}(\xi)$ as ações de rebalanceamento no cenário $\xi$.

  \item \textbf{Otimização Robusta sob Ambiguidade}
    \citep{BenTalNemirovski2002,Fabozzi2007}: especificar um conjunto
    de ambiguidade $\mathcal{U}$ para os parâmetros Nelson-Siegel e
    resolver o minimax:
    \begin{equation}
      \max_{\mathbf{x}} \min_{\xi \in \mathcal{U}}
        \text{NIM}(\mathbf{x}, \xi)
      \label{eq:robust}
    \end{equation}

  \item \textbf{Função Objetivo Multi-Dimensional} com Pareto-ótimo:
    \begin{equation}
      \max_{\mathbf{x}} \Bigl[
        \text{NIM}(\mathbf{x}),\;
        -\Delta\text{EVE}(\mathbf{x}),\;
        \text{RAROC}_{\text{total}}(\mathbf{x})
      \Bigr]
      \label{eq:multiobj}
    \end{equation}

  \item \textbf{Controle Estocástico / MDP}
    \citep{BertsekasSheve1978,Costa2005}: modelar o rebalanceamento
    periódico como Processo de Decisão de Markov com estado
    $s_t = (\text{balanço}_t, \mathbf{r}_t, \text{spreads}_t)$,
    ação $a_t = \mathbf{x}_t$ e recompensa $r_t = \text{NIM}_t$.
\end{enumerate}

\subsection{Posicionamento da Contribuição}

\begin{table}[H]
\centering
\caption{Posicionamento do Modelo Proposto vs.\ Literatura}
\label{tab:posicionamento}
\begin{tabular}{lcccc}
\toprule
\textbf{Dimensão} &
  \textbf{Klein-Monti} &
  \textbf{Kouw.\,\&\,Zenios} &
  \textbf{Baseline QP} &
  \textbf{Proposto} \\
\midrule
Horizonte           & Estático   & Multi-per.  & Um per.   & Multi-per. \\
Incerteza taxa      & Nenhuma    & Cenários    & Paramét.  & Estoc./Robusta \\
Liquidez regulat.   & Não        & Parcial     & Completa  & Completa+din. \\
Capital econômico   & Não        & Não         & RAROC IRB & RAROC + EVE \\
Crédito direcionado & Não        & Não         & Não       & Explícito \\
Solver              & Analítico  & SLP         & SLSQP     & Estocástico \\
\bottomrule
\end{tabular}
\fonte{Elaborado pelo autor, com base em \citet{Klein1971},
       \citet{KouwenbergZenios2001} e \citet{BirgeLouveaux2011}.}
\end{table}

%% =============================================================
\section{Arquitetura do Sistema Computacional}
\label{sec:arquitetura}
%% =============================================================

\subsection{Estrutura de Módulos}

O sistema é implementado em Python~3.12, com separação clara entre dados,
modelos e saídas:

\begin{lstlisting}[language=bash,
                   caption={Estrutura de diretórios do projeto ALM Bank},
                   label={lst:estrutura}]
alm_bank/
|-- main.py                        # Orquestrador (--fase 1..5)
|-- models/
|   |-- balance_sheet/
|   |   |-- balance_sheet.py       # M1: Balanco patrimonial
|   |   `-- market_scenarios.py    # M2: Curvas Nelson-Siegel
|   |-- cash_flow/
|   |   `-- cash_flow_engine.py    # M3: NII, duration, DV01
|   |-- optimization/
|   |   |-- alm_optimizer.py       # M4: QP (SLSQP)
|   |   `-- market_environment.py  # M4a: Mercado disponivel
|   |-- liquidity/
|   |   `-- liquidity_module.py    # M6: LCR, NSFR, stress test
|   |-- capital/
|   |   `-- capital_module.py      # M7: RAROC, CE, pricing
|   `-- dashboard/
|       `-- dashboard_generator.py # M8: Dashboard HTML
|-- data/
|   |-- processed/                 # ativos.csv, passivos.csv, ...
|   `-- scenarios/                 # curvas_di.csv, spreads, ...
|-- outputs/
|   |-- charts/                    # 20 figuras PNG
|   |-- tables/                    # CSV/JSON por modulo
|   `-- reports/                   # dashboard_alm.html
`-- docs/                          # Este documento LaTeX
\end{lstlisting}

\subsection{Fluxo de Execução}

\begin{description}
  \item[\textbf{Fase 1}] Balanço com 35 produtos, calibração Nelson-Siegel,
    3 cenários. Saídas: \texttt{ativos.csv}, \texttt{kpis\_balanco.json}.
  \item[\textbf{Fase 2}] Motor de fluxo de caixa: 10 \textit{buckets},
    NII/NIM mensais, \textit{duration}/DV01 por produto.
    Saídas: \texttt{gap\_reprecificacao.csv}, \texttt{duration\_gap.json}.
  \item[\textbf{Fase 3}] Otimizador QP, 28 variáveis, 5 restrições,
    fronteira eficiente com 12 pontos.
    Saídas: \texttt{fronteira\_eficiente\_base.csv}.
  \item[\textbf{Fase 4}] Stress test LCR, NSFR, custo HQLA, RAROC, pricing.
    Saídas: \texttt{stress\_test\_lcr.csv}, \texttt{raroc\_linhas\_negocio.csv}.
  \item[\textbf{Fase 5}] Dashboard HTML com 6 abas:
    \texttt{outputs/reports/dashboard\_alm.html}.
\end{description}

%% =============================================================
\section{Conclusões}
\label{sec:conclusoes}
%% =============================================================

Este relatório apresentou a implementação completa de um sistema de ALM
bancário em Python, calibrado para o mercado brasileiro de dezembro de 2024,
que constitui o \textit{baseline} quantitativo da pesquisa de doutorado no
PPGEE-Poli/USP. Os principais resultados para o Banco Modelo S.A.\ são:

\begin{enumerate}
  \item \textbf{NIM base}: R\$\,61,0\,bi/ano, NIM de 11,47\%; estável entre
    cenários ($\pm 0{,}3$\,bi), refletindo alto grau de indexação ao CDI.

  \item \textbf{Risco de taxa}: \textit{Duration gap} de 1,742 anos
    (\textit{long}), DV01 gap de R\$\,92,7\,M/bp. A fronteira eficiente
    demonstra que hedgear completamente custa apenas 0,54\% do NIM.

  \item \textbf{Liquidez}: LCR de 1.311\% e NSFR de 160,9\%; o NSFR é a
    restrição ativa no otimizador. Custo do buffer HQLA: \brlM{9.400}/ano.

  \item \textbf{Capital econômico}: RAROC total de 159\%; quatro segmentos
    abaixo do \textit{hurdle rate} de 15\% \citep{Bonomo2016}.

  \item \textbf{Pricing}: Cinco segmentos com \textit{spread} atual acima
    do mínimo; quatro abaixo do mínimo para RAROC $\geq 15\%$.
\end{enumerate}

As limitações identificadas --- dinâmica multi-período do NSFR, assimetria
EVE vs.\ NIM e tratamento do crédito direcionado --- fornecem motivação
clara para a tese:
\textit{Otimização Estocástica Bancária do NIM para Guidance ALM-ALCO
sob Contornos Regulatórios em Mercados Completos}.

%% =============================================================
%%  REFERÊNCIAS
%%  ABNT NBR 6023:2018
%%  natbib + hyperref: cada \citep{} e \citet{} no texto
%%  gera um link clicável (cor azul) para a entrada abaixo.
%%  Não é necessário rodar BibTeX — thebibliography é manual.
%% =============================================================
\newpage
\renewcommand{\bibname}{REFERÊNCIAS}

%% bibitem format for natbib authoryear:
%%   \bibitem[Label(Ano)]{chave}
%%   onde Label é o que aparece nas citações: (AUTOR, ANO)
\begin{thebibliography}{99}

\bibitem[BCBS(2006)]{BCBS2006}
  BASEL COMMITTEE ON BANKING SUPERVISION.
  \textbf{International Convergence of Capital Measurement and Capital
  Standards: A Revised Framework (Basel II)}.
  Basel: Bank for International Settlements, 2006.

\bibitem[BCBS(2013)]{BCBS238}
  BASEL COMMITTEE ON BANKING SUPERVISION.
  \textbf{Basel III: The Liquidity Coverage Ratio and Liquidity Risk
  Monitoring Tools}. BCBS 238.
  Basel: Bank for International Settlements, 2013.

\bibitem[BCBS(2014)]{BCBS295}
  BASEL COMMITTEE ON BANKING SUPERVISION.
  \textbf{Basel III: The Net Stable Funding Ratio}. BCBS 295.
  Basel: Bank for International Settlements, 2014.

\bibitem[BCBS(2016)]{BCBS368}
  BASEL COMMITTEE ON BANKING SUPERVISION.
  \textbf{Interest Rate Risk in the Banking Book}. BCBS 368.
  Basel: Bank for International Settlements, 2016.

\bibitem[Ben-Tal e Nemirovski(2002)]{BenTalNemirovski2002}
  BEN-TAL, A.; NEMIROVSKI, A.
  Robust optimization: methodology and applications.
  \textbf{Mathematical Programming}, Heidelberg, v.~92, n.~3,
  p.~453--480, 2002.

\bibitem[Bertsekas e Shreve(1978)]{BertsekasSheve1978}
  BERTSEKAS, D.~P.; SHREVE, S.~E.
  \textbf{Stochastic Optimal Control: The Discrete Time Case}.
  New York: Academic Press, 1978.

\bibitem[Birge e Louveaux(2011)]{BirgeLouveaux2011}
  BIRGE, J.~R.; LOUVEAUX, F.
  \textbf{Introduction to Stochastic Programming}.
  2.~ed. New York: Springer, 2011.

\bibitem[Bonomo e Martins(2016)]{Bonomo2016}
  BONOMO, M.; MARTINS, B.
  The impact of government-driven loans in the monetary transmission
  mechanism: what can we learn from firm-level data?
  \textbf{Banco Central do Brasil Working Paper}, Brasília, n.~419, 2016.

\bibitem[Coelho et al.(2012)]{Coelho2012}
  COELHO, C.~A.; DE MELLO, J.~M.~P.; FUNCHAL, B.
  The Brazilian payroll lending experiment.
  \textbf{Review of Economics and Statistics}, Cambridge, v.~94, n.~4,
  p.~925--934, 2012.

\bibitem[Costa et al.(2005)]{Costa2005}
  COSTA, O.~L.~V.; FRAGOSO, M.~D.; MARQUES, R.~P.
  \textbf{Discrete-Time Markov Jump Linear Systems}.
  London: Springer, 2005.

\bibitem[Drechsler et al.(2017)]{Drechsler2017}
  DRECHSLER, I.; SAVOV, A.; SCHNABL, P.
  The deposits channel of monetary policy.
  \textbf{The Quarterly Journal of Economics}, Oxford, v.~132, n.~4,
  p.~1819--1876, 2017.

\bibitem[Fabozzi(1996)]{Fabozzi1996}
  FABOZZI, F.~J.
  \textbf{Fixed Income Mathematics}.
  3.~ed. Chicago: Probus Publishing, 1996.

\bibitem[Fabozzi et al.(2007)]{Fabozzi2007}
  FABOZZI, F.~J.; FOCARDI, S.~M.; KOLM, P.~N.
  \textbf{Robust Portfolio Optimization and Management}.
  Hoboken: Wiley, 2007.

\bibitem[Klein(1971)]{Klein1971}
  KLEIN, M.~A.
  A theory of the banking firm.
  \textbf{Journal of Money, Credit and Banking}, Columbus, v.~3, n.~2,
  p.~205--218, 1971.

\bibitem[Koussis e Padgett(2016)]{KoussisPadgett2016}
  KOUSSIS, N.; PADGETT, C.
  Optimal bank capital structure with NSFR constraints.
  \textbf{Journal of Financial Stability}, Amsterdam, v.~25,
  p.~123--135, 2016.

\bibitem[Kouwenberg e Zenios(2001)]{KouwenbergZenios2001}
  KOUWENBERG, R.; ZENIOS, S.~A.
  Stochastic programming models for asset-liability management.
  In: ZENIOS, S.~A.; ZIEMBA, W.~T. (ed.).
  \textbf{Handbook of Asset and Liability Management}.
  Amsterdam: Elsevier, 2001. v.~1.

\bibitem[Kraft(1988)]{Kraft1988}
  KRAFT, D.
  A software package for sequential quadratic programming.
  \textbf{Technical Report DFVLR-FB 88-28}.
  Köln: DLR German Aerospace Center, 1988.

\bibitem[Kusy e Ziemba(1986)]{KusyZiemba1986}
  KUSY, M.~I.; ZIEMBA, W.~T.
  A bank asset and liability management model.
  \textbf{Operations Research}, Catonsville, v.~34, n.~3,
  p.~356--376, 1986.

\bibitem[Merton e Perold(1993)]{MertonPerold1993}
  MERTON, R.~C.; PEROLD, A.~F.
  Theory of risk capital in financial firms.
  \textbf{Journal of Applied Corporate Finance}, Hoboken, v.~6, n.~3,
  p.~16--32, 1993.

\bibitem[Monti(1972)]{Monti1972}
  MONTI, M.
  Deposit, credit and interest rate determination under alternative bank
  objectives.
  In: SZEGO, G.~P.; SHELL, K. (ed.).
  \textbf{Mathematical Methods in Investment and Finance}.
  Amsterdam: North-Holland, 1972.

\bibitem[Nelson e Siegel(1987)]{NelsonSiegel1987}
  NELSON, C.~R.; SIEGEL, A.~F.
  Parsimonious modeling of yield curves.
  \textbf{The Journal of Business}, Chicago, v.~60, n.~4,
  p.~473--489, 1987.

\bibitem[Nocedal e Wright(2006)]{Nocedal2006}
  NOCEDAL, J.; WRIGHT, S.~J.
  \textbf{Numerical Optimization}.
  2.~ed. New York: Springer, 2006.

\bibitem[Resti e Sironi(2007)]{ResteSironi2007}
  RESTI, A.; SIRONI, A.
  \textbf{Risk Management and Shareholders' Value in Banking}.
  Chichester: Wiley, 2007.

\bibitem[Saunders e Walter(2011)]{SaundersWalter2011}
  SAUNDERS, A.; WALTER, I.
  \textbf{Financial Architecture, Systemic Risk, and Universal Banking}.
  Cambridge: MIT Press, 2011.

\end{thebibliography}

\end{document}
